\section{Two further producers: inverse matrices and multivariate charts}
\label{sec:applications}

The matrix machinery has uses beyond iterating a single update.
This section gives two explicit producers. The first constructs an
inverse name from a finite contraction certificate. The second
constructs a passing multivariate chart certificate from a regular
rational base point. Both use finite identities, with all precision
choices determined before execution.

\subsection{A finite Neumann inverse}

For positive weights \(w\), define the induced matrix bound
\[
 \norm{M}_w=\max_{i<d}\frac{\sum_{j<d}|M_{ij}|w_j}{w_i}.
\]
The triangle inequality is immediate from finite sums. Expanding a
matrix product gives
\begin{align*}
 \sum_j|(MN)_{ij}|w_j
 &\le\sum_t|M_{it}|\sum_j|N_{tj}|w_j\\
 &\le\norm{N}_w\sum_t|M_{it}|w_t,
\end{align*}
so \(\norm{MN}_w\le\norm M_w\norm N_w\). All sums are finite,
and reordering them is a finite distributive identity. Also
\(\norm I_w=1\).

Let \(K\) be a rational matrix with \(|K|\le A\), where
\(A,w,s\) pass the small-gain test and \(q=1-1/s\). Then
\(\norm K_w\le q\). Define
\[
 S_N=\sum_{j=0}^{N-1}K^j,\qquad S_0=0.
\]

\begin{theorem}[Constructed inverse near the identity]
\label{thm:inverse}
The finite matrices \(S_N\) satisfy
\begin{align}
 (I-K)S_N=S_N(I-K)&=I-K^N,                     \tag{M1}\\
 \norm{S_M-S_N}_w&\le s q^N\quad(M\ge N),    \tag{M2}\\
 \norm{S_N}_w&\le s.                          \tag{M3}
\end{align}
Choose \(H\) with \(s\le2^H\), and set
\[
 C(k)=S_{s(k+H+2)}.
\]
This is a PR regular matrix sampler. It represents a two-sided inverse
of \(I-K\), with an explicit residual rate.
\end{theorem}
\begin{proof}
Expand the products in (M1). Every intermediate power cancels in
the finite telescoping sum. Since \(K\) commutes with its own
powers, the same identity holds on both sides.
For (M2), submultiplicativity gives
\[
 \norm{S_M-S_N}_w\le\sum_{j=N}^{M-1}q^j
 =q^N\sum_{j=0}^{M-N-1}q^j\le s q^N.
\]
The last inequality follows without an infinite series: multiply
the finite sum by \(1-q=1/s\) to get \(1-q^{M-N}\le1\).
This also handles \(q=0\), with \(q^0=1\). Taking \(N=0\)
gives (M3). The block bound \(q^s\le1/2\) gives
\(s q^{s(k+H+2)}\le2^{-k-2}\), so (M2) proves regularity.
Finally (M1) bounds both inverse residuals by \(q^{s(k+H+2)}\).
The all-precision equations for a two-sided inverse therefore follow
from explicit finite inequalities.
\end{proof}

This matrix norm controls every entry:
\(\abs{M_{ij}}\le(w_i/w_j)\norm M_w\). To return entry
\((i,j)\) with the conventional dyadic error, shift the matrix
precision by a boundedly computed binary magnitude of \(w_i/w_j\).
Thus the matrix-valued notation entails an ordinary family of
rational coordinate names.

For a general rational matrix \(D\), a rational preconditioner
\(B\) can be used by setting \(K=I-BD\). If a rational
two-sided inverse \(B^{-1}\) is also supplied and checked, then
\(D=B^{-1}(I-K)\), and the constructed inverse is \(C B\).
At a finite cutoff, the two residuals are explicitly
\[
 (S_NB)D-I=-K^N,\qquad D(S_NB)-I=-B^{-1}K^NB.
\]
Their norm bounds are respectively \(q^N\) and
\(\norm{B^{-1}}_w\norm B_w q^N\). A binary magnitude of
\(\max(1,\norm{B^{-1}}_w\norm B_w)\) supplies the additional
precision shift for the right residual. Multiplication by \(B\)
in the regular sampler similarly uses a shift for \(\norm B_w\).
This formulation deliberately makes the preconditioner's invertibility
evidence visible. One may replace it with a finite determinant or
elimination proof, but should then include that proof rather than
silently cancel a possibly singular matrix.

\begin{proposition}[Stability of the inverse construction]
Suppose \(K,L\) both have weighted norm at most the same \(q<1\).
Their inverse names differ, in represented matrix norm, by at most
\(s^2\norm{K-L}_w\), where \(s=(1-q)^{-1}\) is an integer
as in the preceding construction.
\end{proposition}
\begin{proof}
For every positive integer \(j\), finite telescoping gives
\[
 K^j-L^j=\sum_{t=0}^{j-1}K^t(K-L)L^{j-1-t}.
\]
Summing through \(N-1\), the norm of the difference of the two
partial Neumann sums is at most
\[
 \norm{K-L}_w\sum_{a+b\le N-2}q^{a+b}
 \le\norm{K-L}_w
       \left(\sum_{a<N}q^a\right)
       \left(\sum_{b<N}q^b\right)
 \le s^2\norm{K-L}_w.
\]
The cases \(N=0,1\) are empty sums. Comparing different cutoffs
adds only the explicit tails (M2). At any requested precision,
choose the two cutoffs to make those tails smaller than the assigned
slack. This is the represented inequality asserted.
\end{proof}

\subsection{Automatic multivariate local charts at rational base points}

Let \(u\in\Q^d\) be unknown coordinates and \(v\in\Q^e\)
parameters, with \(d\ge1\). Let \(P(v,x)\in\Q[v,x]^d\) be
a finite polynomial system, and let \(c\in\Q^d\) satisfy the
coefficient-checkable equality \(P(0,c)=0\). Write
\(J=\partial_xP(0,c)\), computed by formal differentiation of
the polynomial table. Suppose rational matrices \(B,C\) pass
\[
 BJ=I,\qquad CB=I.                             \tag{M4}
\]
For example, when a nonzero rational determinant of \(J\) has
been checked, one can use its finite adjugate formula for \(B\)
and set \(C=J\). The identities are finite rational tests.

Expand the explicit polynomial
\[
 T(v,u)=u-BP(v,c+u)=
       \sum_{\alpha,\beta}t_{i,\alpha,\beta}u^\alpha v^\beta
       \quad\text{in coordinate }i.           \tag{M5}
\]
Equation (M4) and the base-point equality imply
\[
 t_{i,0,0}=0,\qquad t_{i,e_j,0}=0
 \quad(i,j<d).
\]
These are proved coefficient identities, not zero tests on
arbitrary represented values.

Compute nonnegative rational constants
\begin{align*}
 C_0&=\max_i\sum_{\alpha=0,\,|\beta|\ge1}
                       |t_{i,\alpha,\beta}|,\\
 C_1&=\max_i\sum_{|\alpha|=1,\,|\beta|\ge1}
                       |t_{i,\alpha,\beta}|,\\
 C_2&=\max_i\sum_{|\alpha|\ge2}
                       |\alpha|\,|t_{i,\alpha,\beta}|.
\end{align*}
All omitted sums are zero. Now set
\[
 r=\min\left(1,\frac1{4\max(1,C_2)}\right),
 \qquad
 h=\min\left(1,\frac{r}{2\max(1,C_0)},
                   \frac1{4\max(1,C_1)}\right).
 \label{eq:chart-radii}\tag{M6}
\]
Both radii are explicitly positive rationals.

\begin{theorem}[Produced multivariate chart certificate]
\label{thm:chart}
For every \(v\) with \(\norm v_\infty\le h\), the map
\(u\mapsto T(v,u)\) preserves \(\norm u_\infty\le r\) and
contracts by \(1/2\). Consequently it has a uniformly PR regular
solution sampler \(u(v,k)\) in that box, and
\(x(v,k)=c+u(v,k)\) represents a root of \(P(v,x)=0\).
This root is unique in \(\norm{x-c}_\infty\le r\).
The construction is uniform in the finite polynomial code, dimensions,
base point, and inverse certificate.
\end{theorem}
\begin{proof}
On the parameter box, the constant-in-\(u\) part of (M5) has
norm at most \(hC_0\le r/2\). For the difference between two
\(u\) inputs, factor each monomial difference coordinate by
coordinate. The sum of the absolute coefficients of the resulting
Lipschitz bound is at most
\[
 \sum_{\alpha,\beta}|\alpha|\,
 |t_{i,\alpha,\beta}|r^{|\alpha|-1}h^{|\beta|}.
\]
Terms with \(|\alpha|=0\) do not occur in this sum. The pure
linear terms vanish. The remaining linear terms are at most
\(hC_1\le1/4\); all terms with \(|\alpha|\ge2\) are at
most \(rC_2\le1/4\), because \(r,h\le1\).
More explicitly, take the componentwise gain entries
\[
 A_{ij}=\sum_{\alpha,\beta:\,|\alpha|\ge1}
 \alpha_j |t_{i,\alpha,\beta}|
 r^{|\alpha|-1}h^{|\beta|}.
\]
Monomial factorization gives the componentwise action inequality for
this matrix, and the row sums just proved are at most \(1/2\).
Thus \(A\one\le\one/2\), the literal small-gain certificate.
Comparing an input with zero proves
\(\norm{T(v,u)}_\infty\le r/2+\norm u_\infty/2\le r\).

The exact polynomial action, regarded as an approximate action
independent of its precision argument, satisfies the coupled
producer hypotheses with \(w=\one\), \(s=2\), and diameter
\(2r\). Apply \cref{thm:small-gain-producer}. Its finite
fixed-point residual is an error bound on \(BP(v,c+u)\).
Multiplication by the supplied \(C\) in (M4) gives an explicit
error bound on \(P(v,c+u)\) itself. For example, if
\(M=\max_i\sum_j|C_{ij}|\), a shift by a binary magnitude
of \(2M\) makes every coordinate residual at most the requested
dyadic error. All these operations are bounded rational operations.

For uniqueness, any represented root in the box is a represented
fixed point of (M5). Comparing its finite observations with the
constructed sampler uses the residual separation inequality (SG6).
The residuals can be made as small as desired by their explicit
moduli, yielding equality. No new existence assertion about
completed reals is needed.
\end{proof}

\subsection{Parameter dependence and named inputs}

Finite factorization in the parameter coordinates gives the computable
bound
\[
 K_v=\max_i\sum_{\alpha,\beta:,|\beta|\ge1}
   |\beta|\,|t_{i,\alpha,\beta}|
              r^{|\alpha|}h^{|\beta|-1}
\]
for the sup-norm Lipschitz constant in \(v\), on the certified
boxes. The perturbation theorem (SG16) therefore gives a solution
modulus with constant \(2K_v\). In elementary notation,
\[
 \norm{u(v)-u(v')}_\infty\le2K_v\norm{v-v'}_\infty.
\]
Its precise meaning is the corresponding inequality between
finite samples with their known dyadic slack.

If a parameter is supplied by a registered regular name, clip its
rational samples coordinatewise to the parameter box. Clipping
does not increase coordinate distance. Choose the working
parameter precision large enough that \(K_v\) times the input
error is at most the desired action error. Then the polynomial
coefficient construction in \cref{sec:small-gain} proves (SG3)
for the resulting rational evaluator. The same producer thus
handles those named parameters as well.

The theorem constructs local multivariate solution charts; a finite
collection of such charts can be compared using their explicit
uniqueness and parameter bounds. It does not infer that such charts
cover an arbitrary algebraic set, and it does not handle a singular
Jacobian without additional mathematical structure.
